基因编辑效率的机器学习预测已取得进展，但预测精度的提升并未同步带来对影响因素的可靠解释：线性模型方向明确但表达受限，树模型的重要性指标缺乏方向，深度模型的内部表示难以直接解读。针对这一"预测—解释"缺口，本研究构建了一个基因编辑影响因素科学发现平台，把质量控制、特征工程、受控多模型实验、模型归因、统计推断、环境因素析因分析、序列模式挖掘、细胞背景比较与跨模型证据整合组织为一条可追溯的分析流程，并以 DeepCRISPR 公开数据（4 个细胞系、16\,749 条 23\,nt 靶序列、1\,344 次受控训练运行）开展案例研究。平台使用 Linear Regression、XGBoost、多层感知机（MLP）、双分支卷积神经网络（CNN，序列卷积核 3/5/7）与 Transformer 五类具有不同归纳偏置的模型，并严格区分预测效应（$\Delta R^{2}$）、模型归因、稳健性与统计证据。

案例研究的主要观察包括：（i）五类模型均取得正的中位留出集 $R^{2}$（0.035--0.115），说明数据中存在可学习但较弱的编辑效率信号；（ii）四个模型类（CNN、MLP、Transformer、XGBoost）独立地将 PAM 邻近种子区（第 17--20 位）识别为最高归因区域，第 18 位在跨模型平均归因谱中排序最高，且该位置 C 相对 A 的实测效率差在 HCT116（$+0.090$）与 HeLa（$+0.092$）中较大、在 HL60 中减小（$+0.028$）、在 HEK293T 中方向反转（$-0.015$），提示该位置的碱基关联具有细胞背景依赖性；（iii）在 192 组严格配对实验中，CNN 序列卷积核 5 与 7 相对核 3 的 $\Delta R^{2}$ 分别为 $+0.042$ 与 $+0.056$，且归因向 PAM 邻近区域集中，而三个配置提取的候选模式长度中位数均为 4\,nt，说明卷积核是模型感受野参数而非模式长度；（iv）环境因素的跨模型平均增量预测价值很小（$|\Delta R^{2}|<0.01$），2\,541 个边$\times$种子检验中仅 257 个置信区间不跨 0，区组析因方差分析未检出显著主效应（$p=0.056$--$0.663$）；（v）从 CNN 归因中提取到 596 个候选序列模式，其中 64 个通过富集错误发现率校正，但这些候选模式目前仅具计算证据。

平台的科学性不体现在"所有分析均得到阳性结果"，而体现在统一的、可追溯的框架中区分稳定证据、模型特异性现象、上下文依赖现象与不确定结果，并据此输出可检验的候选假设（如第 18 位碱基替换、PAM 邻近候选模式扰动）。本研究属于**计算案例研究（computational case study）**：已完成受控计算验证、统计推断与候选因素优先级排序，尚未开展湿实验验证；平台输出的是待验证假设，而非已确认的生物学机制。
